\begin{thebibliography}{99}

\bibitem{bhr:ProveIt}
V.~Reshetnikov, \emph{ProveIt: polylogarithm manuscript and incoming
research collection}, repository snapshot
\texttt{7bd45777a8a127e5dc69aff8887ca36a79a3059d},
11 October 2026.
\href{https://github.com/VladimirReshetnikov/ProveIt/tree/7bd45777a8a127e5dc69aff8887ca36a79a3059d/Analysis/Polylogarithms/docs/manuscript}{Pinned manuscript};
\href{https://github.com/VladimirReshetnikov/ProveIt/tree/7bd45777a8a127e5dc69aff8887ca36a79a3059d/docs/incoming}{pinned incoming collection}.
The attached inventory lists all source files used for the targeted audit.

\bibitem{bhr:Parity}
ProveIt research collection,
\emph{Harmonic Parity and Resolvent Identities},
incoming report dated 11 October 2026, in \cite{bhr:ProveIt}.
Sections 5 and 7; in particular the local remainder and symmetric
germ with labels \src{qtr:polar}, \src{qtr:kernel}, \src{qtr:J},
and questions R4, R6, R10.
\href{https://github.com/VladimirReshetnikov/ProveIt/blob/7bd45777a8a127e5dc69aff8887ca36a79a3059d/docs/incoming/ProveIt_Harmonic_Parity_and_Resolvent_Identities_2026-10-11.zip}{Pinned source archive}.

\bibitem{bhr:Cubic}
ProveIt research collection,
\emph{Cubic Tornheim Derivatives and Harmonic Stieltjes Constants},
incoming report dated 11 October 2026, in \cite{bhr:ProveIt}.
The cubic digamma--harmonic bridge and Section 7, question 6.
\href{https://github.com/VladimirReshetnikov/ProveIt/blob/7bd45777a8a127e5dc69aff8887ca36a79a3059d/docs/incoming/ProveIt_Cubic_Harmonic_Mixed_Orders_2026-10-11.zip}{Pinned source archive}.

\bibitem{bhr:DLMFZ}
NIST Digital Library of Mathematical Functions,
\S25.11, \emph{Hurwitz Zeta Function}.
\url{https://dlmf.nist.gov/25.11}.
Fourier expansion, argument shifts, Bernoulli values, and derivatives.
Accessed 11 October 2026.

\bibitem{bhr:DLMFG}
NIST Digital Library of Mathematical Functions,
\S\S5.7, 5.11, 5.15, 5.17, \emph{Gamma Function}.
\url{https://dlmf.nist.gov/5.15};
\url{https://dlmf.nist.gov/5.17}.
Digamma and polygamma expansions, asymptotics, and the Barnes $G$ function.
Accessed 11 October 2026.

\bibitem{bhr:DLMFPolylog}
NIST Digital Library of Mathematical Functions,
\S25.12, equation 25.12.12, \emph{Polylogarithms}.
\url{https://dlmf.nist.gov/25.12.E12}.
The Jonqui\`ere expansion used in Mellin subtractions.
Accessed 11 October 2026.

\bibitem{bhr:EMIntegrals}
O.~Espinosa and V.~H.~Moll,
\emph{On some integrals involving the Hurwitz zeta function: Part 2},
Ramanujan Journal \textbf{6} (2002), 449--468.
\url{https://www.math.tulane.edu/~vhm/papers_html/hurwitz2.pdf}.

\bibitem{bhr:EMGeneralized}
O.~Espinosa and V.~H.~Moll,
\emph{A generalized polygamma function},
Integral Transforms and Special Functions \textbf{15} (2004), 101--115.
\url{https://arxiv.org/abs/math/0305079}.

\bibitem{bhr:Adamchik}
V.~S.~Adamchik,
\emph{Multiple Gamma Function and Its Application to Computation of Series},
arXiv:math/0308074 (2003).
\url{https://arxiv.org/abs/math/0308074}.
Classical multiple-Gamma normalization and Hurwitz-zeta conversion.

\bibitem{bhr:Vardi}
I.~Vardi,
\emph{Determinants of Laplacians and multiple gamma functions},
SIAM Journal on Mathematical Analysis \textbf{19} (1988), 493--507.
\url{https://doi.org/10.1137/0519035}.

\bibitem{bhr:Mezo}
I.~Mez\H{o},
\emph{The Fourier series of the log-Barnes function},
arXiv:1604.00753v1 (2016).
\url{https://arxiv.org/pdf/1604.00753v1}.
Known Fourier coefficients; the pointwise correction in
Section~\ref{bhr:sec:audit} concerns equation (7), printed page 6 of this
specific version, and preserves its integrated use.

\bibitem{bhr:Kargin}
L.~Karg\i n, A.~Dil, M.~Cenkci, and M.~Can,
\emph{On the Stieltjes constants with respect to harmonic zeta functions},
Journal of Mathematical Analysis and Applications \textbf{525} (2023),
article 127302.
\url{https://arxiv.org/abs/2304.03517};
\url{https://doi.org/10.1016/j.jmaa.2023.127302}.
The harmonic-zeta coefficient framework is prior work; our strict
depth-three normalization is specified independently in the text.

\bibitem{bhr:Nakamura}
T.~Nakamura,
\emph{Rapidly convergent series representations of symmetric Tornheim
double zeta functions}, arXiv:2103.16873 (2021), Theorem 1.2.
\url{https://arxiv.org/abs/2103.16873}.

\end{thebibliography}
